\section{Síntesis y ampliación}

Esta sesión parte de la pregunta «qué es una representación visual general» y, pasando por VTAB, BiT, ViT, Scaling ViT y MLP-Mixer, construye una metodología más duradera que memorizar arquitecturas:

\begin{enumerate}
  \item \textbf{Definir primero el objetivo de adaptación}: una representación general debe reducir el costo en muestras, en optimización y en especificación de las tareas nuevas;
  \item \textbf{Medir con familias de tareas y no con una sola tabla de clasificación}: las natural, specialized y structured tasks revelan dimensiones distintas de la generalización;
  \item \textbf{Distinguir upstream de downstream}: la representación preentrenada y la capacidad del adaptador deben controlarse por separado;
  \item \textbf{Sin etiquetas no significa sin conocimiento previo}: el SSL traslada el diseño al proxy objective y a la augmentation;
  \item \textbf{La escala es una variable conjunta}: los datos, la capacidad, la duración del entrenamiento y la regularization deben corresponderse;
  \item \textbf{La deduplicación es un requisito de validez científica}: los casi duplicados hacen pasar la memorización por transferencia;
  \item \textbf{ViT es un experimento que debilita los conocimientos previos visuales}: el patch, la position y el class token siguen constituyendo decisiones de diseño explícitas;
  \item \textbf{Los datos determinan el intervalo en el que el sesgo inductivo aporta valor}: un conocimiento previo fuerte mejora la eficiencia con pocos datos, y un modelo flexible puede alcanzarlo con muchos datos;
  \item \textbf{Compute-normalized sigue sin equivaler a eficiencia del sistema}: los FLOPs, la latency, la memory y la utilization deben reportarse por separado;
  \item \textbf{La shape y el schedule forman parte del scaling}: un modelo grande no puede evaluarse con el régimen de entrenamiento corto de un modelo pequeño;
  \item \textbf{La scaling law es un modelo empírico local}: sirve para planificar, pero hay que reportar el intervalo de ajuste y el riesgo de puntos de inflexión;
  \item \textbf{El sentido de Mixer es una comparación contrafactual}: la attention no es el único mecanismo; el objeto causal es la recipe completa.
\end{enumerate}

\begin{importantbox}{Tarea del curso: diseñar una evaluación confiable de representaciones visuales}
Elige un modelo preentrenado y tres downstream tasks claramente distintas. Define la upstream provenance, el método de deduplicación y dos protocol, linear probe y full fine-tuning; reporta para cada grupo de tareas la media y el peor caso (mean/worst-case), 10-shot/full-data, la calibration, los FLOPs de training/inference y la latency real. Por último, escribe al menos tres «cosas que los resultados no pueden demostrar» y explica, si el presupuesto de datos se multiplicara por diez, en qué familia de modelos esperarías que cambiara la pendiente de la curva.
\end{importantbox}

\begin{warningbox}{Autoevaluación final: no mezclar cinco tipos de éxito}
\textbf{Upstream fit}, \textbf{in-distribution accuracy}, \textbf{few-shot transfer}, \textbf{robustness} y \textbf{deployment efficiency} son cinco tipos distintos de éxito. Si un artículo o un sistema solo demuestra uno de ellos, la conclusión debe limitarse a ese; no se puede sustituir la representación general por el top-1 de ImageNet, ni el costo real de servicio por los FLOPs teóricos.
\end{warningbox}

\subsection{Lecturas adicionales}

{\small
\begin{itemize}[nosep,leftmargin=*]
  \item \href{https://arxiv.org/abs/1910.04867}{\textbf{VTAB}}: benchmark de adaptación de representaciones visuales multitarea y diseño por task-family.
  \item \href{https://arxiv.org/abs/1912.11370}{\textbf{Big Transfer (BiT)}}: preentrenamiento supervisado a gran escala, recipe de transferencia y robustness.
  \item \href{https://arxiv.org/abs/2010.11929}{\textbf{An Image is Worth 16x16 Words}}: arquitectura ViT, scaling de datos y transfer.
  \item \href{https://arxiv.org/abs/2106.10270}{\textbf{How to Train Your ViT?}}: augmentation, regularization, dataset size y training recipe.
  \item \href{https://arxiv.org/abs/2106.04560}{\textbf{Scaling Vision Transformers}}: model shape, schedule, few-shot y scaling laws.
  \item \href{https://arxiv.org/abs/2105.01601}{\textbf{MLP-Mixer}}: token/channel mixing y data-scale comparison.
\end{itemize}
}

\end{document}
